\begin{tikzpicture}[
    scale=0.85,
    line width=1pt,
    minimum size=0.08cm
]

% 1 rdzeń
% =====================================================================
\begin{scope}[shift={(0,0)}]

    \node[circle, draw=black, fill=cfyel!85]  (v1) at (0,2) {};
    \node[circle, draw=black, fill=cfblue!85] (v2) at (2,2) {};
    \node[circle, draw=black, fill=cfyel!85]  (v3) at (2,0) {};
    \node[circle, draw=black, fill=cfblue!85] (v4) at (0,0) {};
    \node[circle, draw=black, fill=cfred!85]  (v5) at (1,1) {};

    \draw (v1) -- (v2) -- (v3) -- (v4) -- (v1);

    \draw (v2) -- (v5);
    \draw (v3) -- (v5);

    \node at (1,-0.8) {1 kopia rdzenia};

\end{scope}


% kilka kopii
% =====================================================================
\begin{scope}[shift={(5,0)}]

    \node[circle, draw=black, fill=cfyel!85]  (v1) at (0,2) {};
    \node[circle, draw=black, fill=cfblue!85] (v2) at (2,2) {};
    \node[circle, draw=black, fill=cfyel!85]  (v3) at (2,0) {};
    \node[circle, draw=black, fill=cfblue!85] (v4) at (0,0) {};
    \node[circle, draw=black, fill=cfred!85]  (v5) at (1,1) {};

    \draw (v1) -- (v2) -- (v3) -- (v4) -- (v1);

    \draw (v2) -- (v5);
    \draw (v3) -- (v5);
    \draw (v4) -- (v5);

    \node at (1.2,-0.8) {kilka kopii rdzenia};

\end{scope}


% kopie pokrywają cały graf
% =====================================================================
\begin{scope}[shift={(10,0)}]

    \node[circle, draw=black, fill=cfyel!85]  (v1) at (0,2) {};
    \node[circle, draw=black, fill=cfblue!85] (v2) at (2,2) {};
    \node[circle, draw=black, fill=cfyel!85]  (v3) at (2,0) {};
    \node[circle, draw=black, fill=cfblue!85] (v4) at (0,0) {};
    \node[circle, draw=black, fill=cfred!85]  (v5) at (1,1) {};

    \draw (v1) -- (v2) -- (v3) -- (v4) -- (v1);

    \draw (v1) -- (v5);
    \draw (v2) -- (v5);
    \draw (v3) -- (v5);
    \draw (v4) -- (v5);


    \node at (1,-0.8) {kopie rdzenia};
    \node at (1,-1.4) {pokrywają cały graf};

\end{scope}

\end{tikzpicture}